\documentclass[aspectratio=169]{beamer}

\usetheme{Madrid}
\usecolortheme{default}
\usefonttheme{professionalfonts}

\usepackage{amsmath,amssymb,bm}
\usepackage{siunitx}
\usepackage{physics}
\usepackage{mathtools}

% Vectors
\newcommand{\grad}{\bm{\nabla}}
\newcommand{\J}{\bm{J}}
\newcommand{\N}{\bm{N}}
\newcommand{\jvec}{\bm{j}}
\newcommand{\uvec}{\bm{u}}

\title[Fick for Gases]{Engineering Transport Derivation of Fick's Law for Gases}
\subtitle{From Species Conservation + Stefan--Maxwell to Mixture-Averaged / Dilute Limits}
\author{}
\institute{}
\date{}

\begin{document}

% ==========================================================
\begin{frame}
  \titlepage
\end{frame}

% ==========================================================
\begin{frame}{What ``Fick's law'' means in gas-phase engineering}
In gases, engineers distinguish:
\begin{itemize}
  \item \textbf{Total species flux} (in a fixed lab frame): $\N_i$ [\si{mol.m^{-2}.s^{-1}}]
  \item \textbf{Diffusive flux relative to the mixture motion}: $\J_i$ or $\jvec_i$
  \item \textbf{Bulk (convective) transport} carried by the mixture velocity $\uvec$
\end{itemize}

\vspace{0.3cm}
You will see \emph{multiple} ``Fick forms'' depending on variables and frame:
\[
\J_i = - c\,D_{i,m}\,\grad x_i \quad \text{or} \quad
\jvec_i = - \rho\,D_{i,m}\,\grad Y_i
\]
where $x_i$ = mole fraction, $Y_i$ = mass fraction, $c$ = total molar concentration, $\rho$ = density.
\end{frame}

% ==========================================================
\begin{frame}{Step 1: Species conservation (continuum, reacting/nonreacting)}
Molar form for species $i$:
\[
\boxed{\pdv{c_i}{t} + \grad\cdot \N_i = \dot{\omega}_i}
\]
\begin{itemize}
  \item $c_i$ [\si{mol.m^{-3}}], $\N_i$ [\si{mol.m^{-2}.s^{-1}}]
  \item $\dot{\omega}_i$ is molar production by reactions (set to 0 for nonreacting)
\end{itemize}

Mass form:
\[
\boxed{\pdv{(\rho Y_i)}{t} + \grad\cdot(\rho Y_i \uvec + \jvec_i) = \dot{\omega}_{i,\text{mass}}}
\]
Key: we must model $\N_i$ (or $\jvec_i$).
\end{frame}

% ==========================================================
\begin{frame}{Step 2: Flux decomposition and frames}
Define mixture-average velocity $\uvec$ (mass-average is common in CFD):
\[
\uvec \equiv \sum_i Y_i \uvec_i
\]

Total molar flux in the lab frame decomposes as:
\[
\boxed{\N_i = c_i \uvec + \J_i}
\]
where
\begin{itemize}
  \item $c=\sum_i c_i$ (total molar concentration),
  \item $\J_i$ is the \textbf{diffusive molar flux relative to $\uvec$}.
\end{itemize}

Mass-flux analogue:
\[
\boxed{\rho Y_i \uvec + \jvec_i}
\qquad\text{with}\qquad \sum_i \jvec_i = \bm{0}
\]
(when $\uvec$ is mass-average).
\end{frame}

% ==========================================================
\begin{frame}{Step 3: What drives diffusion in gases (engineering view)}
In gas mixtures, diffusion is driven by gradients in:
\begin{itemize}
  \item composition ($x_i$, $Y_i$),
  \item pressure (barodiffusion),
  \item temperature (thermal diffusion / Soret effect),
  \item external body forces (e.g., electric fields for ions).
\end{itemize}

\vspace{0.2cm}
\textbf{Engineering simplification (most common):}
\begin{itemize}
  \item ideal gas mixture
  \item low Mach / modest pressure gradients
  \item negligible Soret and barodiffusion
  \item isothermal or weak $\grad T$
\end{itemize}
Then diffusion is mainly due to \textbf{composition gradients}.

\vspace{0.25cm}
The rigorous multicomponent closure is Stefan--Maxwell.
\end{frame}

% ==========================================================
\begin{frame}{Step 4: Stefan--Maxwell equations (multicomponent gases)}
For an ideal gas mixture (no body forces, ignoring Soret/barodiffusion), Stefan--Maxwell gives:
\[
\boxed{
\grad x_i = \sum_{j\neq i}\frac{x_i x_j}{D_{ij}}\left(\frac{\uvec_j-\uvec_i}{1}\right)\frac{p}{p}
}
\]

Engineering form in terms of \textbf{diffusive molar fluxes} relative to a chosen reference velocity:
\[
\boxed{
-\grad x_i = \sum_{j\neq i}\frac{x_j \J_i - x_i \J_j}{c\,D_{ij}}
}
\]
where
\begin{itemize}
  \item $D_{ij}$ is the \textbf{binary (Maxwell--Stefan) diffusivity} of pair $(i,j)$,
  \item $c=p/(RT)$ for ideal gases.
\end{itemize}

This is the ``first-principles'' engineering starting point for gas diffusion closures.
\end{frame}

% ==========================================================
\begin{frame}{Constraint: fluxes must sum to zero in the mixture frame}
If $\uvec$ is the \textbf{molar-average} velocity, then
\[
\sum_i \J_i = \bm{0}
\]
If $\uvec$ is the \textbf{mass-average} velocity, then
\[
\sum_i \jvec_i = \bm{0}
\]

These constraints matter because \textbf{Stefan--Maxwell is coupled}: $\J_i$ depends on all $\J_j$.
Fick’s law emerges only after a simplifying limit/closure.
\end{frame}

% ==========================================================
\begin{frame}{Binary mixture: Stefan--Maxwell $\Rightarrow$ Fick exactly}
Consider only species $A$ and $B$.

Stefan--Maxwell reduces to:
\[
-\grad x_A = \frac{x_B \J_A - x_A \J_B}{c\,D_{AB}}
\]
With the molar-average frame constraint $\J_A+\J_B=0 \Rightarrow \J_B=-\J_A$:
\[
-\grad x_A = \frac{(x_B+x_A)\J_A}{c\,D_{AB}} = \frac{\J_A}{c\,D_{AB}}
\]
Hence:
\[
\boxed{\J_A = -c\,D_{AB}\,\grad x_A}
\]
This is a clean derivation: \textbf{in a binary ideal gas}, Fick’s law holds (for the diffusive flux in the molar-average frame).
\end{frame}

% ==========================================================
\begin{frame}{From diffusive flux to total flux (what you measure)}
Total molar flux:
\[
\N_A = c_A \uvec + \J_A
\]
Thus, even with Fickian diffusion:
\[
\boxed{\N_A = c_A \uvec - c\,D_{AB}\,\grad x_A}
\]

Two important special cases:
\begin{itemize}
  \item \textbf{No bulk motion} ($\uvec=\bm{0}$): $\N_A = -cD_{AB}\grad x_A$
  \item \textbf{With bulk flow}: diffusion rides on top of convection
\end{itemize}

This is why CFD codes always write:
\[
\text{(convection)} + \text{(diffusion)}
\]
\end{frame}

% ==========================================================
\begin{frame}{Multicomponent mixture: mixture-averaged approximation}
For $n$ species, solving full Stefan--Maxwell is expensive.

A common engineering closure: \textbf{mixture-averaged diffusion}
\[
\boxed{\J_i \approx -c\,D_{i,m}\,\grad x_i}
\]
where $D_{i,m}$ is the effective diffusivity of species $i$ into the mixture.

A widely used formula (good for ideal gases, moderate gradients) is:
\[
\boxed{
D_{i,m} = \frac{1-x_i}{\sum\limits_{j\neq i}\frac{x_j}{D_{ij}}}
}
\]

Notes:
\begin{itemize}
  \item If species $i$ is dilute ($x_i\ll 1$), then $D_{i,m}\approx \left(\sum_{j\neq i}x_j/D_{ij}\right)^{-1}$.
  \item Coupling is collapsed into a single scalar $D_{i,m}$.
\end{itemize}
\end{frame}

% ==========================================================
\begin{frame}{Dilute-species limit: Fick becomes natural}
Let species $A$ be dilute in a carrier gas mixture (everything except $A$ is the ``bath''):
\[
x_A \ll 1,\quad x_{\text{bath}}\approx 1
\]
Then (mixture-averaged):
\[
\J_A \approx -c\,D_{A,m}\,\grad x_A,
\qquad
D_{A,m}\approx \left(\sum_{j\neq A}\frac{x_j}{D_{Aj}}\right)^{-1}
\]
If the bath is essentially a \textbf{single} carrier $B$:
\[
D_{A,m}\approx D_{AB}
\]
So you recover the familiar binary Fick form with a physically meaningful diffusivity.
\end{frame}

% ==========================================================
\begin{frame}{Engineering forms: mole fraction vs concentration vs partial pressure}
For ideal gases:
\[
c=\frac{p}{RT},\qquad c_i = x_i c,\qquad p_i = x_i p
\]

From
\[
\J_i = -c D_{i,m}\grad x_i
\]
you can write equivalent versions:

\textbf{In terms of concentration:}
\[
\J_i = -D_{i,m}\left(\grad c_i - x_i \grad c\right)
\]
\textbf{At constant $p,T$ (so $c$ constant):}
\[
\boxed{\J_i = -D_{i,m}\,\grad c_i}
\]
\textbf{In terms of partial pressure (constant $T$):}
\[
\boxed{\J_i = -\frac{D_{i,m}}{RT}\,\grad p_i}
\]
These are the standard engineering identities used in gases.
\end{frame}

% ==========================================================
\begin{frame}{Mass-fraction form (common in CFD)}
Many solvers use mass fractions $Y_i$ and mass-average velocity.

A typical mixture-averaged closure is:
\[
\boxed{\jvec_i \approx -\rho\,D_{i,m}\,\grad Y_i}
\]
with a correction (in some implementations) to enforce $\sum_i \jvec_i=\bm{0}$ exactly:
\[
\jvec_i = -\rho D_{i,m}\grad Y_i + Y_i\sum_k \rho D_{k,m}\grad Y_k
\]
So the \textbf{engineering meaning of “Fick’s law”} often is:
\[
\text{diffusive mass flux} \propto -\grad Y_i
\]
with a mixture constraint correction.
\end{frame}

% ==========================================================
\begin{frame}{From flux law + species conservation to diffusion equation}
Nonreacting, no bulk flow, constant $p,T$ so $c$ constant:

Species conservation:
\[
\pdv{c_i}{t} + \grad\cdot \J_i = 0
\]
Fick/mixture-averaged:
\[
\J_i = -D_{i,m}\grad c_i
\]
Thus:
\[
\boxed{\pdv{c_i}{t} = \grad\cdot\left(D_{i,m}\grad c_i\right)}
\]
If $D_{i,m}$ is constant:
\[
\boxed{\pdv{c_i}{t} = D_{i,m}\nabla^2 c_i}
\]
This is the gas-phase diffusion equation in its simplest engineering form.
\end{frame}

% ==========================================================
\begin{frame}{Important caveat: diffusion in gases can induce bulk flow}
In gas mixtures, diffusion may create a net molar flux unless constrained.

Classic example: \textbf{Stefan diffusion} (evaporation of $A$ through stagnant $B$).
Even if $B$ is ``stagnant'', $A$ diffusion can drive a bulk velocity to satisfy overall continuity and boundary constraints.

Engineering takeaway:
\begin{itemize}
  \item Use $\N_i = c_i \uvec + \J_i$, not just $\N_i=\J_i$
  \item ``Stagnant'' does not always mean $\uvec=\bm{0}$
  \item Stefan--Maxwell captures this naturally; pure Fick needs care
\end{itemize}
\end{frame}

% ==========================================================
\begin{frame}{How gas diffusivity depends on $T$ and $p$ (engineering scaling)}
For dilute gases, binary diffusivities roughly scale as:
\[
D_{AB} \propto \frac{T^{3/2}}{p}
\]
More detailed kinetic-theory correlations exist (Chapman--Enskog / collision integrals),
but the key engineering trend is:

\begin{itemize}
  \item higher $T$ $\Rightarrow$ faster molecular motion $\Rightarrow$ larger $D$
  \item higher $p$ (higher density) $\Rightarrow$ more collisions $\Rightarrow$ smaller $D$
\end{itemize}

\vspace{0.2cm}
When you see tables of $D_{AB}$, check the stated $T$ and $p$.
\end{frame}

% ==========================================================
\begin{frame}{Summary: what we derived and when it is valid}
\textbf{Starting point (multicomponent, gas-phase):}
\[
-\grad x_i = \sum_{j\neq i}\frac{x_j \J_i - x_i \J_j}{cD_{ij}}
\quad\text{(Stefan--Maxwell)}
\]

\textbf{Binary ideal gas:} Fick is exact for diffusive flux:
\[
\J_A = -cD_{AB}\grad x_A
\]

\textbf{Multicomponent engineering closure:} mixture-averaged diffusivity:
\[
\J_i \approx -cD_{i,m}\grad x_i,\qquad
D_{i,m} = \frac{1-x_i}{\sum_{j\neq i}x_j/D_{ij}}
\]

\textbf{Always remember:} total flux = convection + diffusion
\[
\N_i = c_i\uvec + \J_i
\]
\end{frame}

% ==========================================================
\begin{frame}{Optional slide: where to go beyond Fick for gases}
If you need higher fidelity:
\begin{itemize}
  \item include Soret (thermal diffusion): $\grad T$ drives $\J_i$
  \item include barodiffusion: $\grad p$ effects
  \item use full multicomponent Stefan--Maxwell solver
  \item consider non-ideal mixtures (real-gas, high pressure): use chemical potentials/fugacities
  \item include Knudsen diffusion (porous media, small pores)
\end{itemize}

\vspace{0.25cm}
In combustion and high-temperature flows, mixture-averaged is common; multicomponent can matter for light species (H, H$_2$).
\end{frame}

\end{document}
